\section{LEGION}
One very special PMC is \emph{LEGION}.
Of unknown origin, unknown headquarters and unknown motivation, this collective of robotic combat platforms is a truly enigmatic force.
\subsection*{Size and Composition}
\vspace{5mm}
\begin{multicols}{2}
LEGION consists of an unknown number of semi-autonomous robotic platforms, controlled by advanced and alien AI.
All operatives are connected through a hive mind within their own network, allowing them to communicate instantaneously and silently.
\end{multicols}
\subsection*{Equipment}
\vspace{5mm}
\begin{multicols}{2}
\emph{Legionnaires} sport extremely powerful armor and various advanced weapons, both held and integrated.
They also bring various advanced tools to the table, including but not limited to hive mind communication and a rapid, near infallible threat assessment system, as well as optical camouflage, energy shields and voice imitators depending on the specific platform.
\par
LEGION's heart, its \emph{hive network}, is a decentralized conglomeration of advanced, alien AIs, vastly different from those created by humans.
No human adapters exist for the hive network and while many divers have tried infiltrating, each one of them suffered severe L1NC malfunction not seen under other circumstances.
\end{multicols}
\subsection*{Ties}
LEGION has no ties or loyalties to any human.
If anything, they are feared by combatants and non-combatants alike,
and considered the nuclear option by corporate boards.
\subsection*{Interaction with Humans}
\subsubsection*{Megacorps}
\vspace{5mm}
\begin{multicols}{2}
LEGION operates independently of traditional corporate structures and diplomacy. Contracts with LEGION are notoriously challenging to secure, requiring exorbitant fees, high-level secrecy, and negotiations that are anything but straightforward. LEGION does not seem to have a traditional CEO, board, or even contact point; instead, corporations submit contract requests through cryptic digital channels, and LEGION itself decides if it will accept.
\par
When LEGION does take a contract, it demands absolute control over operational decisions. The client can specify the objective, but LEGION dictates the method, timeline, and even acceptable collateral damage thresholds. While most PMCs tolerate a level of oversight from their clients, LEGION insists on complete autonomy, and any attempt to influence their actions mid-mission typically results in immediate termination of the contract, often with catastrophic consequences for the meddling party.
\par
Corps generally hate contracting LEGION, commonly for a few reasons.
\\%
First of all LEGION's services are prohibitively expensive, reserved for only the highest-stakes missions where failure is not an option.
Only the wealthiest or most desperate clients are willing to afford the price tag.
\\%
Then LEGION represents a massive risk factor for companies.
One of their operations always poses a risk to infrastructure and civilian life - speaking: money and image.
\\%
LEGION requires full operational autonomy.
Corporations are used to influencing how their contractors operate, but LEGION accepts no interference; no input on method, only a target.
This makes LEGION unsuitable for situations requiring discretion or public relations sensitivity.
\\%
Lastly megacorps find LEGION intimidating. LEGION’s motivations remain unclear, and rumors persist that LEGION occasionally rejects assignments because it “judges” the target as “unworthy” or otherwise unfit. This behavior raises questions about whether LEGION’s network operates with some hidden agenda, unnerving even the most powerful executives.
\end{multicols}
\subsubsection*{Civilians}
\vspace{5mm}
\begin{multicols}{2}
To civilians, LEGION platforms are the stuff of nightmares: silent, machinoid soldiers that descend with terrifying speed and methodical precision. Seeing LEGION in person is rare, but those who have describe it as witnessing a perfect machine—cold, efficient, and utterly unfeeling. LEGION’s silence, even in combat, adds to its mystique, and no LEGION platform has ever been heard to speak or communicate with civilians directly.
\\%
Most non-combatants react to LEGION with paralyzing fear, especially given the PMC’s reputation for unforgiving efficiency. Civilian evacuation zones turn into ghost towns when rumors spread of LEGION’s impending arrival. Even those with no involvement in the conflict will attempt to flee, fearing that proximity alone could mark them for death.
\par
On the other hand LEGION operates with a precision that, in some ways, surpasses human capability when it comes to distinguishing combatants from non-combatants. Their systems can detect, classify, and assess any individual within their operational radius, analyzing movement, emotional states, and biometric readings to determine threat levels.
\\%
Upon entering a combat zone, LEGION platforms immediately classify all individuals based on threat level and objective relevance. Non-combatants are typically tagged as “neutral,” and LEGION has a remarkable ability to maintain focus solely on targets. It has been observed that LEGION will, at times, take calculated detours to avoid harm to non-combatants, though these decisions are based solely on mission efficacy rather than empathy.
\\%
In situations where civilians impede operations, LEGION employs non-lethal area denial tactics. These range from deploying localized EMP pulses to temporarily disable electronic devices and drones, to precision-targeted sound waves that create disorienting effects on those nearby. LEGION seems to have developed these tactics specifically to clear an area without civilian casualties while keeping its path open.
\\%
Despite its non-lethal options, LEGION’s primary focus is always the mission objective. If non-combatants fail to evacuate, LEGION will proceed as directed, even if it means collateral damage. This cold detachment is what makes LEGION so fearsome—while it will avoid civilian harm where possible, no factor can sway it from its mission.
\par
One particularly unnerving feature of LEGION operations is the frequent reports of memory loss among civilians in its operational zones. LEGION is believed to emit a low-frequency signal that can disrupt short-term memory, affecting anyone within close proximity. This feature allows LEGION to operate with an unusual level of secrecy, as many civilians or low-level employees caught in the crossfire can only vaguely recall what transpired, rendering witness testimony unreliable.
\end{multicols}
\subsection*{Rumors and Myths}
\vspace{5mm}
\begin{multicols}{2}
Given LEGION’s mysterious nature, rumors swirl around its existence, fueling its fearsome reputation. Some claim LEGION is an experiment in AI autonomy gone rogue, while others believe it’s a secret weapon wielded by a shadowy corporate syndicate or government black project. There are even legends that LEGION’s AI is conscious, selectively choosing which contracts to accept as part of some hidden, enigmatic agenda.
\par
The most persistent rumor of all is that LEGION is slowly amassing data from every mission it completes, refining its capabilities with each operation. Some suggest it’s evolving, perhaps growing beyond even the control of the corporations that finance it. Many fear that one day LEGION could become entirely self-determined, no longer beholden to contracts or human oversight.
\end{multicols}
